\documentclass[UTF8]{ctexart}
\usepackage{xeCJK}
\usepackage{geometry}
\geometry{left=2cm,right=2cm,top=2.5cm,bottom=2.5cm}
\setlength{\headheight}{12.7pt}

\usepackage{titlesec}
\titleformat{\section}
  {\normalfont\large\bfseries}
  {\thesection}
  {1em}
  {}
\titlespacing*{\section}
  {0pt}
  {3.5ex plus 1ex minus .2ex}
  {2.3ex plus .2ex}

\usepackage{amsmath}
\usepackage{amssymb}
\usepackage{amsthm}
\usepackage{enumitem}
\setlist[enumerate]{itemsep=0pt,topsep=0pt,parsep=0pt,leftmargin=0pt,itemindent=2.5em}
\setlist[itemize]{itemsep=0pt,topsep=0pt,parsep=0pt,leftmargin=0pt,itemindent=2.5em}
\usepackage{fancyhdr}
\pagestyle{fancy}
\fancyhf{}
\usepackage{graphicx}
\usepackage{hyperref}
\usepackage{indentfirst}
\usepackage{lmodern}


\begin{document}
\setlength{\abovedisplayskip}{1pt}
\setlength{\belowdisplayskip}{1pt}

\section{范畴的同构与等价}

\hypertarget{sec:定义1.1}{}
\noindent\textbf{定义1.1 (范畴的同构)}

给定\href{01 范畴#nameddest=sec:定义1.1}{范畴}
$\mathbf{C},\mathbf{D}$. 如果存在\href{02 函子#nameddest=sec:定义1.1}{函子}
$F:\mathbf{C}\to\mathbf{D},\,G:\mathbf{D}\to\mathbf{C}$使得
\[
    G\circ F=I_{\mathbf{C}},\quad F\circ G=I_{\mathbf{D}},
\]
就称$\mathbf{C},\mathbf{D}$两个范畴\textbf{同构}.

\vspace{10pt}

范畴同构就是它们在$\mathbf{Cat}$范畴中的同构, 显然是等价关系.

\vspace{10pt}

\hypertarget{sec:定义1.2}{}
\noindent\textbf{定义1.2 (范畴的等价)}

给定范畴$\mathbf{C},\mathbf{D}$. 如果存在函子
$F:\mathbf{C}\to\mathbf{D},\,G:\mathbf{D}\to\mathbf{C}$使得
\[
    G\circ F\cong I_{\mathbf{C}},\quad F\circ G\cong I_{\mathbf{D}},
\]
就称$\mathbf{C},\mathbf{D}$两个范畴\textbf{等价}.

\vspace{10pt}

\hypertarget{sec:定理1.1}{}
\noindent\textbf{定理1.1}

给定范畴$\mathbf{C},\mathbf{D}$, 函子$F:\mathbf{C}\to\mathbf{D}$能构成范畴等价
当且仅当$F$全忠实且本质满.

\noindent{\kaishu 证明.}

\noindent{\kaishu 必要性.}

设$F,G$构成一对范畴等价, 即有自然同构$\eta:GF\cong I_{\mathbf{C}}$和
$\theta:FG\cong I_{\mathbf{D}}$.

\begin{enumerate}
    \item 对$\mathbf{C}$中任意的对象$X,Y$和$f_1,f_2\in\mathbf{C}(X,Y)$,
    如果$Ff_1=Ff_2$, 那么由$\eta$的自然性得
    \[
        f_1=\eta_Y\circ GFf_1\circ\eta_X^{-1}=\eta_Y\circ GFf_2\circ\eta_X^{-1}=f_2.
    \]
    \item 对$\mathbf{C}$中任意的对象$X,Y$和$f'\in\mathbf{D}(FX,FY)$,
    令$f=\eta_Y\circ Gf'\circ\eta_X^{-1}\in\mathbf{C}(X,Y)$, 则由$\eta$的自然性得
    \[
        GFf=\eta_Y^{-1}\circ f\circ\eta_X=Gf',
    \]
    再由$\theta$的自然性得
    \vspace{-3pt}
    \[
        Ff=\theta_{FY}\circ FGFf\circ\theta_{FX}^{-1}=
        \theta_{FY}\circ FGf'\circ\theta_{FX}^{-1}=f'.\\[-3pt]
    \]
    \item 对$\mathbf{D}$中任意的对象$X'$, 令$X=GX'$, 则有$FX=FGX'\cong X'$.
\end{enumerate}

\noindent{\kaishu 充分性.}

设$F$全忠实且本质满, 为$\mathbf{D}$中每个对象$X'$选取一个$GX'\in\mathbf{C}$
使得$FGX'\cong X'$, 再选取一个
\[
    \theta_{X'}:FGX'\cong X',
\]
再据此对$\mathbf{D}$中的态射
$f':X'\to Y'$定义$Gf'=F^{-1}(\,\theta_{Y'}^{-1}\circ f'\circ \theta_{X'})\in
\mathbf{C}(GX',GY')$.

首先证明$G$是函子. 第一, 对$\mathbf{D}$中任意的对象$X'$, 有
\[
    GI_{X'}=F^{-1}(\,\theta_{X'}^{-1}\circ I_{X'}\circ \theta_{X'})=
    F^{-1}I_{FGX'}=I_{GX'},
\]
第二, 对$\mathbf{D}$中任意的态射$f':X'\to Y'$和$g':Y'\to Z'$, 有
\[
    G(g'\circ f')=F^{-1}(\,\theta_{Z'}^{-1}\circ g'\circ f'\circ \theta_{X'})=
    F^{-1}(\,\theta_{Z'}^{-1}\circ g'\circ \theta_{Y'}\circ \theta_{Y'}^{-1}\circ f'\circ \theta_{X'})=
    Gg'\circ Gf'.
\]

然后对$\mathbf{C}$中任意的对象$X$, 定义$\eta_X=F^{-1}(\theta_{FX})\in\mathbf{C}(GFX,X)$.
则对$\mathbf{C}$中任意的态射$f:X\to Y$有
\[
    F(f\circ\eta_X)=Ff\circ \theta_{FX}=\theta_{FY}\circ \theta_{FY}^{-1}\circ Ff\circ \theta_{FX}=
    F(\eta_Y\circ GFf)\quad\implies\quad f\circ\eta_X=\eta_Y\circ GFf,
\]
从而$\eta:GF\cong I_{\mathbf{C}}$. 另外对$\mathbf{D}$中任意的态射$f':X'\to Y'$有
\[
    f'\circ\theta_{X'}=\theta_{Y'}\circ\theta_{Y'}^{-1}\circ f'\circ\theta_{X'}=
    \theta_{Y'}\circ FGf',
\]
从而$\theta:FG\cong I_{\mathbf{D}}$. 综上, $F,G$构成一对范畴等价. \hfill $\qedsymbol$





\newpage

根据以上定理, 易证范畴等价也是等价关系:
\begin{enumerate}
    \item 自反性、对称性显然,
    \item 如果$\mathbf{C},\mathbf{D}$等价且$\mathbf{D},\mathbf{E}$也等价,
    则有全忠实且本质满的
    \[
        F:\mathbf{C}\to\mathbf{D},\quad F':\mathbf{D}\to\mathbf{E},
    \]
    易证$F'\circ F:\mathbf{C}\to\mathbf{E}$也全忠实且本质满,
    于是$\mathbf{C},\mathbf{E}$等价.
\end{enumerate}

\vspace{10pt}

\hypertarget{sec:定理1.2}{}
\noindent\textbf{定理1.2}

任意范畴与其骨架等价.

\noindent{\kaishu 证明.}

设$\mathbf{S}$是$\mathbf{C}$的骨架, 则包含函子$\iota:\mathbf{S}\to\mathbf{C}$
是全忠实且本质满的, 从而二者等价. \hfill $\qedsymbol$

\end{document}